\documentclass[11pt]{article}
\usepackage[margin=1in]{geometry}
\usepackage{hyperref}
\usepackage{enumitem}
\usepackage{booktabs}
\usepackage{xcolor}
\title{ProbOS --- Month 3, Week 9\\Extend REST API Model Registry}
\author{Nisong Monyimba}
\date{\today}

\begin{document}
\maketitle

\section*{Week 9 goal}
Extend \texttt{python/server/main.py}'s model registry beyond
\texttt{BatteryModel2Cell} to the three Week 4 cross-discipline models
(\texttt{OptionPricerModel}, \texttt{EDQueueModel},
\texttt{ClinicalTrialModel}) -- closing one of the open SBIR Phase I
asks (\texttt{docs/sbir\_phase1\_onepager.md}: ``Extend the REST API's
model registry beyond \texttt{BatteryModel2Cell} to the
finance/operations/clinical models already built''). This is
deliberately low-risk: the endpoint pattern is already proven for one
model, and the three target models are unchanged, already-validated
Week 4 code.

\section*{Day-by-day plan}

\subsection*{Day 1 -- Registry design + schema audit}
\begin{itemize}[leftmargin=*]
  \item Review each of the three models' constructors and
    prior-builder functions to determine what request-schema fields
    each needs (e.g. \texttt{OptionPricerModel} needs S0/K/T/r;
    \texttt{EDQueueModel} needs Q0; \texttt{ClinicalTrialModel} needs
    randomisation ratio)
  \item Decide whether \texttt{SimulateRequest} gets model-specific
    optional fields, or whether each model gets its own request
    schema class -- lean toward the latter for type safety, following
    the existing \texttt{schemas.py} pattern
  \item Update \texttt{\_build\_model\_and\_priors()}'s registry dict
    to include all four models
\end{itemize}

\subsection*{Day 2 -- /simulate for the three new models}
\begin{itemize}[leftmargin=*]
  \item Extend \texttt{POST /simulate} to accept
    \texttt{model\_name} values \texttt{"option\_pricer"},
    \texttt{"ed\_queue"}, \texttt{"clinical\_trial"}
  \item Verify each returns percentiles/convergence matching a direct
    Python \texttt{MonteCarloEngine} call, at the same
    \texttt{rtol=1e-10} standard already established for
    \texttt{BatteryModel2Cell}
\end{itemize}

\subsection*{Day 3 -- /sensitivity for the three new models}
\begin{itemize}[leftmargin=*]
  \item Extend \texttt{POST /sensitivity} similarly
  \item Confirm \texttt{SobolSensitivity} runs correctly against each
    model's own \texttt{param\_dim}/\texttt{state\_dim} (these differ
    across the four models -- e.g. \texttt{ClinicalTrialModel} has
    \texttt{state\_dim=4}, not 8 like \texttt{BatteryModel2Cell})
\end{itemize}

\subsection*{Day 4 -- /filter for the three new models}
\begin{itemize}[leftmargin=*]
  \item Extend \texttt{POST /filter} to accept the three new models
  \item Confirm the existing Gaussian-observation-noise likelihood
    pattern is appropriate for each, or document where it is not
    (e.g. \texttt{ClinicalTrialModel}'s Week 8 filtering needed a
    fundamentally different, deterministic-state design -- the
    \texttt{/filter} endpoint may need to expose that pattern too, or
    explicitly scope it out for this model with a clear error message
    rather than silently returning a nonsensical result)
\end{itemize}

\subsection*{Day 5 -- Integration tests}
\begin{itemize}[leftmargin=*]
  \item Extend \texttt{python/tests/test\_server.py} with the same
    numerical-match discipline already used for
    \texttt{BatteryModel2Cell} (HTTP response vs direct kernel call at
    \texttt{rtol=1e-10}), for all three new models across all three
    endpoints where applicable
\end{itemize}

\subsection*{Day 6 -- pre\_week\_audit.sh + CI}
\begin{itemize}[leftmargin=*]
  \item Run \texttt{pre\_week\_audit.sh} before considering the week
    complete
  \item No new CI steps expected (existing \texttt{test\_server.py}
    execution already covers this), but verify explicitly
\end{itemize}

\subsection*{Day 7 -- Documentation + retrospective}
\begin{itemize}[leftmargin=*]
  \item Update \texttt{docs/sbir\_phase1\_onepager.md}: move this item
    from the SBIR Ask section to a demonstrated-results bullet, since
    it will now be complete
  \item Update \texttt{README.md}'s REST API section to reflect all
    four registered models
  \item Retrospective appended to this document once complete
\end{itemize}

\section*{Exit criteria}
All three POST endpoints (\texttt{/simulate}, \texttt{/sensitivity},
\texttt{/filter}, where applicable) work correctly for
\texttt{OptionPricerModel}, \texttt{EDQueueModel}, and
\texttt{ClinicalTrialModel}, each verified against a direct Python
kernel call at \texttt{rtol=1e-10} -- the same standard already
established for \texttt{BatteryModel2Cell} in Week 7.

\end{document}
